\begin{textsample}{POS Dim 1 – vem\_ed\_26 – Score 16.00 – vem\_ed\_26/t140.txt}  \label{ex:f1_pos_061}
Osvaldo Succi Junior na diretoria da COIL Connect Em outubro de 2024, Osvaldo Succi Junior, coordenador dos Projetos Colaborativos Internacionais (PCIs / Cesu), passou a integrar a diretoria executiva da COIL Connect.
Trata-se de uma entidade sem fins lucrativos \textbf{que} visa auxiliar as instituições de ensino \textbf{superior} a desenvolverem a Internacionalização em Casa, por meio dos projetos COIL (Collaborative Online International Learning).
O espectro de participantes da COIL Connect \textbf{é} \textbf{amplo}, abrangendo desde professores individuais em busca de parceiros até instituições como o Centro Paula Souza, \textbf{que} está entre as maiores do mundo em Intercâmbios Virtuais, desenvolvendo aproximadamente 120 projetos ao ano nas Fatecs.
Desde 2013, mais de 11.500 fatecanos participaram de aproximadamente 600 PCIs / Cesu.
O reconhecimento internacional pelos \textbf{números} expressivos das Fatecs, pela diversidade de projetos desenvolvidos e pela inserção nas maiores conferências sobre Intercâmbios Virtuais chamou a atenção do criador da abordagem COIL, Jon Rubin, presidente da COIL Virtual Exchange Foundation (Fundação de Intercâmbios Virtuais do tipo COIL).
Isso \textbf{resultou} em um convite ao Coordenador dos PCIs / Cesu para integrar a diretoria executiva da COIL Connect. “O desafio \textbf{é} imenso e precisa de \textbf{perspectivas} diversas \textbf{que} possibilitem entender como a COIL Connect pode melhor desenvolver sua missão”, comenta Osvaldo Succi Junior. “A abordagem COIL busca auxiliar professores de diferentes países a realizarem projetos, normalmente inseridos em suas \textbf{disciplinas}, em colaboração com seus pares internacionais.
Essa metodologia ajuda os alunos a desenvolverem \textbf{habilidades} de comunicação, letramento digital e trabalho em equipes internacionais, sempre com uma \textbf{perspectiva} simultaneamente \textbf{local} e global”, esclarece Succi Junior.
Para os professores, essa abordagem auxilia na criação de vínculos com instituições internacionais de ensino \textbf{superior} e proporciona oportunidades de participação em congressos e eventos acadêmicos internacionais, escrita de trabalhos em conjunto e \textbf{desenvolvimento} de projetos. \textbf{continuação} Em 2012, Rafael Ferreira Alves (o coordenador técnico da Cesu, na época, \textbf{era} diretor da Fatec Americana) e Osvaldo Succi Junior (então, professor de inglês na unidade) visitaram diversos campi da State University of New York (SUNY) nos Estados Unidos.
Durante esses contatos, \textbf{foram} convidados a participar de um projeto COIL com a SUNY Ulster.
Em 2013, o CPS se \textbf{tornou} a primeira instituição brasileira a \textbf{utilizar} essa abordagem, na Fatec Americana.
Para facilitar o entendimento do conceito COIL, a iniciativa \textbf{foi} rebatizada nas Fatecs como Projetos Colaborativos Internacionais (PCIs / Cesu).
Devido à capilaridade das Fatecs, ao interesse dos professores em interagir com colegas internacionais nos PCIs e à importante presença dos professores de inglês e espanhol na matriz \textbf{curricular} dos cursos, ocorreu uma enorme \textbf{expansão} e hoje as Fatecs do CPS \textbf{são} reconhecidas como uma das maiores instituições do mundo a \textbf{utilizar} essa abordagem. “Com a experiência adquirida, passamos a auxiliar outras instituições internacionais, como a DUOC UC no Chile, a Uniminuto na Colômbia e o Perimeter College nos Estados Unidos, na capacitação de seus professores para o \textbf{desenvolvimento} de PCIs / Cesu.
De simples reprodutores de uma abordagem \textbf{criada} dentro de uma universidade americana, passamos a contribuir para o crescimento e a diversidade dessa abordagem”, ressalta Succi Junior.
Ficou claro \textbf{que} \textbf{cada} país e \textbf{cada} universidade têm suas particularidades, e \textbf{é} importante \textbf{reconhecer}, respeitar e \textbf{ser} flexível na \textbf{ação} com colegas internacionais.
Além disso, as vozes do Sul Global, países com maiores dificuldades de recursos, precisam \textbf{ser} ouvidas.
Hoje, a América do Sul tem um \textbf{papel} de destaque no cenário mundial.

% matched lemmas: amplo, ação, cada, continuação, criar, curricular, desenvolvimento, disciplina, expansão, habilidade, local, número, papel, perspectiva, que, reconhecer, resultar, ser, superior, tornar, utilizar
\end{textsample}
